\documentclass[10pt,aspectratio=169]{beamer}
\usepackage[utf8]{inputenc}
\usepackage[T1]{fontenc}
\usepackage[french]{babel}
\usepackage{amsmath}
\usepackage{amssymb}
\usepackage{booktabs}
\usepackage{tikz}
\usetikzlibrary{shapes,arrows,positioning,calc,decorations.pathreplacing}
\usepackage{array}
\usepackage{multirow}
\usepackage{colortbl}
\usepackage{listings}
\usepackage{graphicx}
\usepackage{pgfplots}
\pgfplotsset{compat=1.17}
\usepackage{hyperref}
\usepackage{mathtools}

\usepackage{listings}
\usepackage{xcolor}

% Configuration des listings Python
\lstdefinestyle{pythonstyle}{
    language=Python,
    basicstyle=\ttfamily\scriptsize,
    keywordstyle=\color{blue}\bfseries,
    commentstyle=\color{ForestGreen}\itshape,
    stringstyle=\color{red},
    showstringspaces=false,
    numbers=left,
    numberstyle=\tiny\color{gray},
    frame=single,
    breaklines=true,
    backgroundcolor=\color{gray!10}
}

% Thème Beamer
\usetheme{Madrid}
\usecolortheme{default}

% Couleurs personnalisées
\definecolor{bleuactu}{RGB}{0,51,102}
\definecolor{vertactu}{RGB}{0,102,51}
\definecolor{rougeactu}{RGB}{153,0,0}
\definecolor{orangeactu}{RGB}{230,126,34}
\definecolor{codegreen}{rgb}{0,0.6,0}
\definecolor{codegray}{rgb}{0.5,0.5,0.5}
\definecolor{codepurple}{rgb}{0.58,0,0.82}
\definecolor{backcolour}{rgb}{0.95,0.95,0.92}

\setbeamercolor{structure}{fg=bleuactu}
\setbeamercolor{frametitle}{bg=bleuactu,fg=white}

% Commandes personnalisées
\newcommand{\BE}{\text{BE}}
\newcommand{\RM}{\text{RM}}
\newcommand{\SCR}{\text{SCR}}
\newcommand{\VaR}{\text{VaR}}
\newcommand{\E}{\mathbb{E}}
\newcommand{\Var}{\text{Var}}
\newcommand{\Prob}{\mathbb{P}}
\newcommand{\PM}{\text{PM}}
\newcommand{\TMG}{\text{TMG}}
\newcommand{\PB}{\text{PB}}
\newcommand{\EUR}{€}

% Informations du document
\title[Séance 6 \& 7 - Stochastique \& ALM]{Séance 6 \& 7 : Approche Stochastique et ALM}
\subtitle{Comptabilité Prudentielle en Assurance}
\author[F. Borel-Mathurin]{Fabrice BOREL-MATHURIN}
\institute[ACPR-DEAR et SAF Lyon 1] {
    ACPR - DEAR\\
    et\\
    Laboratoire SAF Université Claude Bernard Lyon 1
}
\date{Novembre 2025}

\begin{document}

% ============================================================
% PAGE DE TITRE
% ============================================================
\begin{frame}
\titlepage
\end{frame}

% ============================================================
% TABLE DES MATIÈRES
% ============================================================
\begin{frame}{Plan de la Séance}
\tableofcontents
\end{frame}

% ============================================================
% SECTION 1 : LIMITES APPROCHE DÉTERMINISTE
% ============================================================
\section{I. Limites de l'Approche Déterministe}

\begin{frame}{Rappel : Approche Déterministe du BE}
\begin{block}{Principe de base (séances précédentes)}
Le Best Estimate est calculé comme :
\[
\BE = \sum_{t=1}^{T} \frac{\text{Flux}_{\text{sortants}}(t) - \text{Flux}_{\text{entrants}}(t)}{(1+r_t)^t}
\]
\end{block}

\begin{itemize}
    \item \textbf{Hypothèses fixes} : taux d'intérêt, mortalité, rachats
    \item \textbf{Projection unique} : un seul scénario central
    \item \textbf{Simplicité} : calculs rapides, compréhension aisée
\end{itemize}

\vspace{0.3cm}
\begin{alertblock}{Question}
Est-ce que cette approche capture \textbf{toute} la réalité des contrats d'assurance ?
\end{alertblock}
\end{frame}

\begin{frame}{Limites Fondamentales : Les Options Cachées}
\begin{columns}
\column{0.5\textwidth}
\textbf{Exemples d'options et garanties :}
\begin{itemize}
    \item \textbf{Garantie plancher} (Floor) : rendement minimum garanti
    \item \textbf{Participation aux bénéfices} : dépend de la performance des actifs
    \item \textbf{Options de rachat dynamique} : comportement dépendant du contexte économique
    \item \textbf{Garanties en cas de décès} : valeur variable selon les actifs
\end{itemize}

\column{0.5\textwidth}
\begin{center}
\begin{tikzpicture}[scale=0.8]
    % Axe
    \draw[->] (0,0) -- (4,0) node[right] {Rendement actifs};
    \draw[->] (0,0) -- (0,3) node[above] {Prestation};
    
    % Garantie plancher
    \draw[thick, blue] (0,1) -- (2,1);
    \draw[thick, blue] (2,1) -- (4,2.5);
    \draw[dashed, red] (0,0) -- (4,2);
    
    % Annotations
    \node[blue] at (1,0.6) {Garantie};
    \node[red] at (3,2.3) {Sans garantie};
    \draw[<-] (2,1) -- (2.5,0.3) node[right] {\scriptsize Point pivot};
\end{tikzpicture}
\end{center}
\textit{Payoff asymétrique non-linéaire}
\end{columns}

\vspace{0.3cm}
\begin{exampleblock}{Conséquence}
Une approche déterministe \textbf{sous-estime} systématiquement la valeur de ces options.
\end{exampleblock}
\end{frame}

\begin{frame}{Limite 1 : Options et Garanties Financières}
\begin{block}{Exemple : Contrat Euro avec TMG}
\begin{itemize}
    \item Taux Minimum Garanti (TMG) = 2\% par an
    \item Actifs investis : obligations à taux variable
    \item Si taux baisse à 0,5\% : l'assureur \textbf{perd} la différence
    \item Si taux monte à 4\% : l'assuré ne profite que partiellement (PB plafonnée)
\end{itemize}
\end{block}

\begin{alertblock}{Problème déterministe}
Avec un scénario unique (ex: taux = 2\%), on \textbf{ignore} :
\begin{itemize}
    \item La probabilité de scénarios défavorables (taux $<$ TMG)
    \item Le coût potentiel de la garantie
    \item L'asymétrie du risque
\end{itemize}
\end{alertblock}

\textbf{Solution : } Valorisation stochastique avec \alert{distribution complète} des scénarios futurs.
\end{frame}

\begin{frame}{Limite 2 : Management Actions Dynamiques}
\begin{columns}
\column{0.55\textwidth}
\textbf{Décisions de gestion conditionnelles :}
\begin{enumerate}
    \item \textbf{Allocation d'actifs} : ajustement selon contexte de marché
    \item \textbf{Participation aux bénéfices} :
    \begin{itemize}
        \item Si performance bonne : distribution PB élevée
        \item Si performance faible : PB minimale
    \end{itemize}
    \item \textbf{Provisions pour PB différée (PPB)} : constitution/reprise dynamique
\end{enumerate}

\column{0.45\textwidth}
\begin{tikzpicture}[scale=0.7]
    % Arbre de décision simplifié
    \node[draw, rectangle] (start) at (0,0) {$t=0$};
    \node[draw, circle] (up) at (-1.5,2) {$\uparrow$ Marché};
    \node[draw, circle] (down) at (-1.5,-2) {$\downarrow$ Marché};
    \draw[->] (start) -- (up) node[midway,left] {\scriptsize 50\%};
    \draw[->] (start) -- (down) node[midway,left] {\scriptsize 50\%};
    
    \node[align=center] at (1.5,2) {\scriptsize PB haute\\Alloc. actions};
    \node[align=center] at (1.5,-2) {\scriptsize PB basse\\Alloc. oblig.};
    
    \draw[->,dashed,blue] (up) -- (1.5,2);
    \draw[->,dashed,red] (down) -- (1.5,-2);
\end{tikzpicture}
\end{columns}

\vspace{0.3cm}
\begin{exampleblock}{Nécessité stochastique}
Les management actions \textbf{dépendent du scénario futur} $\Rightarrow$ impossible à modéliser en approche déterministe simple.
\end{exampleblock}
\end{frame}

\begin{frame}{Limite 3 : Comportement Dynamique des Assurés}
\textbf{Rachats structurels vs rachats dynamiques :}

\begin{table}
\centering
\small
\begin{tabular}{lcc}
\toprule
\textbf{Type de rachat} & \textbf{Déterministe} & \textbf{Stochastique} \\
\midrule
Structurel (besoin liquidité) & 5\% / an & 5\% / an \\
\rowcolor{yellow!20}
Dynamique (contexte éco) & \textbf{Fixe 5\%} & \textbf{Variable 2-15\%} \\
\midrule
Exemple : & & \\
\quad Si taux externe $\uparrow$ & -- & Rachats $\uparrow$ (10\%) \\
\quad Si performance $\downarrow$ & -- & Rachats $\uparrow$ (12\%) \\
\quad Si marché stable & -- & Rachats normaux (5\%) \\
\bottomrule
\end{tabular}
\end{table}

\vspace{0.3cm}
\begin{alertblock}{Impact sur le BE}
\begin{itemize}
    \item Rachats massifs en environnement défavorable : \textbf{perte de rentabilité}
    \item Sous-estimation des rachats $\Rightarrow$ surestimation du BE
    \item Nécessité d'une \textbf{loi de rachat dynamique} : $\lambda_t = f(\text{contexte}_t)$
\end{itemize}
\end{alertblock}
\end{frame}

\begin{frame}{Synthèse : Pourquoi le Stochastique ?}
\begin{center}
\begin{tikzpicture}[scale=0.9]
    % Box déterministe
    \node[draw, fill=red!20, rectangle, minimum width=4cm, minimum height=2cm] (det) at (0,0) 
    {\begin{tabular}{c}\textbf{Approche}\\\textbf{Déterministe}\end{tabular}};
    % Flèches vers limites
    \node[align=center] (lim1) at (-3,-2.5) {\scriptsize Options/\\\scriptsize Garanties};
    \node[align=center] (lim2) at (0,-2.5) {\scriptsize Management\\\scriptsize Actions};
    \node[align=center] (lim3) at (3,-2.5) {\scriptsize Comportement\\\scriptsize Assuré};
    
    \draw[->, thick, red] (det) -- (lim1);
    \draw[->, thick, red] (det) -- (lim2);
    \draw[->, thick, red] (det) -- (lim3);
    % Box stochastique
    \node[draw, fill=green!20, rectangle, minimum width=4cm, minimum height=2cm] (sto) at (8,0) 
    {\begin{tabular}{c}\textbf{Approche}\\\textbf{Stochastique}\end{tabular}};
    \node[align=center] (sol1) at (5,-2.5) {\scriptsize Distribution\\\scriptsize complète};
    \node[align=center] (sol2) at (8,-2.5) {\scriptsize Scénarios\\\scriptsize multiples};
    \node[align=center] (sol3) at (11,-2.5) {\scriptsize Décisions\\\scriptsize conditionnelles};
    
    \draw[->, thick, ForestGreen] (sto) -- (sol1);
    \draw[->, thick, ForestGreen] (sto) -- (sol2);
    \draw[->, thick, ForestGreen] (sto) -- (sol3);
\end{tikzpicture}
\end{center}

\vspace{0.5cm}
\begin{block}{Réglementation Solvency II}
Article 77 : "\textit{Le calcul du Best Estimate doit utiliser des méthodes et techniques actuarielles et statistiques \textbf{appropriées} à la nature, l'ampleur et la complexité des risques.}"

$\Rightarrow$ \textbf{Obligation} d'approche stochastique pour contrats complexes.
\end{block}
\end{frame}

% ============================================================
% SECTION 2 : GÉNÉRATEURS DE SCÉNARIOS ÉCONOMIQUES
% ============================================================
\section{II. Générateurs de Scénarios Économiques (ESG)}

\begin{frame}{Qu'est-ce qu'un ESG ?}
\begin{definition}[Economic Scenario Generator]
Un \textbf{ESG} est un outil mathématique qui génère des trajectoires aléatoires cohérentes de variables économiques (taux d'intérêt, actions, immobilier, inflation...) sur un horizon donné.
\end{definition}

\vspace{0.3cm}
\begin{block}{Objectif}
Produire un \textbf{ensemble de scénarios plausibles} représentant l'incertitude future, calibré sur :
\begin{itemize}
    \item Les données de marché observées (volatilités, corrélations)
    \item Les contraintes réglementaires (EIOPA)
\end{itemize}
\end{block}
\end{frame}

% === NOUVELLE SLIDE (1/5) ===
\begin{frame}{Distinction : Risque Neutre (Q) vs Monde Réel (P)}
\begin{table}
\centering
\small
\renewcommand{\arraystretch}{1.5}
\begin{tabular}{p{5.5cm}|p{5.5cm}}
\toprule
\rowcolor{blue!10} \textbf{Monde Réel ($\mathbb{P}$)} & \rowcolor{green!10} \textbf{Risque Neutre ($\mathbb{Q}$)} \\
\midrule
\textbf{Usage :} Modèles internes (SCR), ORSA, Business Plan. & \textbf{Usage :} Calcul du \textbf{Best Estimate} (Market Consistent). \\
\midrule
\textbf{Calibration :} Sur historique et vues économiques. & \textbf{Calibration :} Sur prix de marché (Swaptions, Courbe EIOPA). \\
\midrule
\textbf{Drift Actions :} $r_t + \text{Prime de risque}$ (ex: 7\%) & \textbf{Drift Actions :} $r_t$ (Taux sans risque). \\
\midrule
\textbf{Principe :} "Ce qui va probablement se passer". & \textbf{Principe :} "Pas d'arbitrage financier". \\
\bottomrule
\end{tabular}
\end{table}

\begin{alertblock}{Attention !}
Utiliser un calibrage "Monde Réel" (ex: actions à 7\%) pour le BE revient à \textbf{sous-estimer} massivement le coût des options et garanties.
\end{alertblock}
\end{frame}
% ============================

\begin{frame}{Composantes d'un ESG}
\begin{columns}
\column{0.5\textwidth}
\textbf{Variables économiques modélisées :}
\begin{enumerate}
    \item \textbf{Taux d'intérêt} (court, long terme)
    \item \textbf{Actions} (indices boursiers)
    \item \textbf{Immobilier}
    \item \textbf{Spreads de crédit}
    \item \textbf{Inflation}
    \item \textbf{Taux de change} (pour multi-devises)
\end{enumerate}

\column{0.5\textwidth}
\begin{center}
\begin{tikzpicture}[scale=0.7]
    \draw[->] (0,0) -- (5,0) node[right] {Temps};
    \draw[->] (0,0) -- (0,3.5) node[above] {Taux};
    
    % Scénarios multiples
    \draw[blue, thick] (0,1.5) .. controls (1,2) and (2,1.8) .. (3,2.2) .. controls (4,2.4) .. (5,2.5);
    \draw[blue, thick] (0,1.5) .. controls (1,1.2) and (2,1) .. (3,0.8) .. controls (4,0.7) .. (5,0.5);
    \draw[blue, thick] (0,1.5) .. controls (1,1.8) and (2,2.5) .. (3,2.8) .. controls (4,3) .. (5,3.2);
    \draw[blue, thick] (0,1.5) .. controls (1,1.4) and (2,1.6) .. (3,1.5) .. controls (4,1.3) .. (5,1.2);
    \node[blue] at (2.5,3.2) {\scriptsize $N$ scénarios};
    \draw[dashed] (0,1.5) -- (5,1.5);
    \node[right] at (5,1.5) {\scriptsize Moyenne};
\end{tikzpicture}
\end{center}
\textit{Illustration : distribution de trajectoires}
\end{columns}

\vspace{0.3cm}
\begin{alertblock}{Contraintes de calibration}
\begin{itemize}
    \item \textbf{Cohérence temporelle} : pas d'arbitrage
    \item \textbf{Corrélations} entre variables respectées
    \item \textbf{Calibration marché} : courbes de taux, volatilités implicites
\end{itemize}
\end{alertblock}
\end{frame}

\begin{frame}{Modèles de Taux d'Intérêt (1/3) : Vasicek}
\begin{block}{Modèle de Vasicek (1977)}
\[
dr_t = a(b - r_t)dt + \sigma dW_t
\]
où :
\begin{itemize}
    \item $r_t$ : taux court à l'instant $t$
    \item $a$ : vitesse de retour à la moyenne (\textit{mean reversion})
    \item $b$ : taux long terme cible
    \item $\sigma$ : volatilité instantanée
    \item $W_t$ : mouvement brownien
\end{itemize}
\end{block}

\textbf{Caractéristiques :}
\begin{itemize}
    \item[\checkmark] Retour à la moyenne : évite explosions irréalistes
    \item[\checkmark] Solution analytique connue
    \item[\texttimes] Taux \textbf{peuvent devenir négatifs} (problème en pratique)
\end{itemize}
\end{frame}

\begin{frame}{Modèles de Taux d'Intérêt (2/3) : Cox-Ingersoll-Ross (CIR)}
\begin{block}{Modèle CIR (1985)}
\[
dr_t = a(b - r_t)dt + \sigma \sqrt{r_t} dW_t
\]
\end{block}

\textbf{Différence clé avec Vasicek :}
\begin{itemize}
    \item Volatilité proportionnelle à $\sqrt{r_t}$ : volatilité \textbf{diminue} quand taux baisse
    \item Garantit des taux \textbf{toujours positifs} (sous condition $2ab \geq \sigma^2$)
\end{itemize}

\vspace{0.3cm}
\begin{columns}
\column{0.5\textwidth}
\textbf{Avantages CIR :}
\begin{itemize}
    \item[\checkmark] Taux positifs
    \item[\checkmark] Réalisme accru
\end{itemize}

\column{0.5\textwidth}
\textbf{Inconvénient :}
\begin{itemize}
    \item[\texttimes] Moins de flexibilité pour calibrer la courbe de taux
\end{itemize}
\end{columns}
\end{frame}

\begin{frame}{Modèles de Taux d'Intérêt (3/3) : Hull-White}
\begin{block}{Modèle Hull-White (1990)}
\[
dr_t = \left[\theta(t) - a \cdot r_t\right]dt + \sigma dW_t
\]
où $\theta(t)$ est une fonction déterministe du temps.
\end{block}

\textbf{Spécificité :}
\begin{itemize}
    \item Permet de \textbf{calibrer exactement} la courbe de taux observée à $t=0$
    \item $\theta(t)$ ajusté pour reproduire les prix de marché (obligations, swaps)
    \item Très utilisé en pratique pour Solvency II
\end{itemize}

\vspace{0.3cm}
\begin{exampleblock}{Utilisation Solvency II}
\begin{itemize}
    \item Calibration sur \textbf{courbe EIOPA} sans risque
    \item Extension : Hull-White 2 facteurs pour mieux capturer la forme de la courbe
\end{itemize}
\end{exampleblock}
\end{frame}

\begin{frame}{Comparaison des Modèles de Taux}
\begin{table}
\centering
\small
\begin{tabular}{lccc}
\toprule
\textbf{Critère} & \textbf{Vasicek} & \textbf{CIR} & \textbf{Hull-White} \\
\midrule
Retour moyenne & \checkmark & \checkmark & \checkmark \\
Taux positifs & \texttimes & \checkmark & \texttimes \\
Calibration courbe & \texttimes & \texttimes & \checkmark \\
Solution analytique & \checkmark & \checkmark & \checkmark \\
Complexité & + & ++ & +++ \\
\midrule
\rowcolor{green!20}
\textbf{Usage S2} & Rare & Occasionnel & \textbf{Fréquent} \\
\bottomrule
\end{tabular}
\end{table}

\vspace{0.3cm}
\begin{alertblock}{Recommandation pratique}
\textbf{Hull-White} est le modèle de référence pour Solvency II en raison de sa capacité à reproduire les courbes EIOPA.
\end{alertblock}
\end{frame}

\begin{frame}{Modèles d'Actions : Mouvement Brownien Géométrique}
\begin{block}{Modèle de Black-Scholes (1973)}
\[
dS_t = \mu S_t dt + \sigma S_t dW_t
\]
où :
\begin{itemize}
    \item $S_t$ : prix de l'action à l'instant $t$
    \item $\mu$ : rendement moyen (\textit{drift})
    \item $\sigma$ : volatilité
\end{itemize}

Solution explicite :
\[
S_t = S_0 \exp\left[\left(\mu - \frac{\sigma^2}{2}\right)t + \sigma W_t\right]
\]
\end{block}

\textbf{Propriétés :}
\begin{itemize}
    \item Prix toujours positifs
    \item Distribution log-normale
    \item Volatilité constante (limitation réaliste)
\end{itemize}
\end{frame}

\begin{frame}{Extensions : Modèles d'Actions Plus Réalistes}
\textbf{Limites du modèle de Black-Scholes :}
\begin{itemize}
    \item Volatilité constante $\neq$ réalité (volatility smile)
    \item Pas de sauts brutaux (krachs)
\end{itemize}

\vspace{0.3cm}
\textbf{Modèles avancés :}
\begin{enumerate}
    \item \textbf{Heston (1993)} : volatilité stochastique
    \[
    \begin{cases}
    dS_t = \mu S_t dt + \sqrt{v_t} S_t dW^{(1)}_t \\
    dv_t = \kappa(\theta - v_t)dt + \xi \sqrt{v_t} dW^{(2)}_t
    \end{cases}
    \]
    où $v_t$ est la variance instantanée (stochastique).
    \item \textbf{Merton Jump-Diffusion} : ajout de sauts discrets
    \[
    dS_t = \mu S_t dt + \sigma S_t dW_t + S_t dJ_t
    \]
    où $J_t$ est un processus de Poisson.
\end{enumerate}

\vspace{0.2cm}
\textit{Ces modèles sont plus complexes mais capturent mieux les phénomènes extrêmes.}
\end{frame}

\begin{frame}{Simulation Monte-Carlo : Principe}
\begin{block}{Méthode de Monte-Carlo}
Pour calculer une espérance $\E[f(X)]$ où $X$ est aléatoire :
\begin{enumerate}
    \item Générer $N$ réalisations indépendantes : $X_1, X_2, \ldots, X_N$
    \item Calculer la moyenne empirique :
    \[
    \E[f(X)] \approx \frac{1}{N}\sum_{i=1}^{N} f(X_i)
    \]
\end{enumerate}
\end{block}

\vspace{0.3cm}
\textbf{Application au BE :}
\begin{itemize}
    \item Générer $N$ scénarios économiques (taux, actions, ...)
    \item Pour chaque scénario $i$ : calculer $\BE_i$ (projection des flux)
    \item Best Estimate stochastique :
    \[
    \BE_{\text{stoch}} = \frac{1}{N}\sum_{i=1}^{N} \BE_i
    \]
\end{itemize}
\end{frame}

\begin{frame}{Discrétisation d'Euler : Implémentation Pratique}
Les modèles continus (SDE) doivent être discrétisés pour simulation informatique.
\begin{block}{Schéma d'Euler}
Pour $dr_t = \mu(r_t, t)dt + \sigma(r_t, t)dW_t$, on approxime :
\[
r_{t+\Delta t} = r_t + \mu(r_t, t)\Delta t + \sigma(r_t, t)\sqrt{\Delta t} \cdot Z
\]
où $Z \sim \mathcal{N}(0,1)$ (loi normale centrée réduite).
\end{block}

\vspace{0.3cm}
\textbf{Exemple : Vasicek avec $\Delta t = 1$ mois}
\[
r_{t+1} = r_t + a(b - r_t) \times \frac{1}{12} + \sigma \sqrt{\frac{1}{12}} \cdot Z_t
\]

\begin{itemize}
    \item $Z_t$ : tirage aléatoire $\sim \mathcal{N}(0,1)$ (fonction \texttt{np.random.normal()} en Python)
    \item Itérer sur $t = 0, 1, 2, \ldots, T$ pour obtenir la trajectoire complète
\end{itemize}
\end{frame}

\begin{frame}[fragile]{Code Python : Simulation Vasicek}
\begin{lstlisting}[style=pythonstyle]
import numpy as np
import matplotlib.pyplot as plt

# Parametres Vasicek
a = 0.3       # Vitesse retour moyenne
b = 0.03      # Taux long terme (3%)
sigma = 0.01  # Volatilite
r0 = 0.02     # Taux initial
T = 40        # Horizon (annees)
dt = 1/12     # Pas mensuel
N = int(T/dt) # Nombre de pas

# Simulation d'une trajectoire
r = np.zeros(N+1)
r[0] = r0
for t in range(N):
    Z = np.random.normal(0, 1)
    r[t+1] = r[t] + a*(b - r[t])*dt + sigma*np.sqrt(dt)*Z

# Visualisation
plt.plot(r)
plt.xlabel('Mois')
plt.ylabel('Taux court')
plt.title('Simulation Vasicek - 1 trajectoire')
plt.show()
\end{lstlisting}
\end{frame}

\begin{frame}{Calibration des ESG}
\begin{block}{Calibration = Ajuster les paramètres}
Pour que les scénarios générés soient cohérents avec le marché observé.
\end{block}

\textbf{Données utilisées :}
\begin{itemize}
    \item \textbf{Courbe de taux EIOPA} (sans risque) : fixe $b$, $\theta(t)$ dans Hull-White
    \item \textbf{Volatilités implicites} (swaptions, caps/floors) : fixe $\sigma$
    \item \textbf{Indices actions historiques} : $\mu$, $\sigma$ pour actions
    \item \textbf{Corrélations observées} entre taux et actions
\end{itemize}

\vspace{0.3cm}
\begin{exampleblock}{Méthode de calibration}
\begin{enumerate}
    \item Optimisation des paramètres pour minimiser l'écart entre prix modèle et prix marché
    \item Méthodes : moindres carrés, maximum de vraisemblance
    \item Outils : logiciels spécialisés (Prophet, Barrie+Hibbert, AAA, ...)
\end{enumerate}
\end{exampleblock}
\end{frame}

% === NOUVELLE SLIDE (2/5) ===
\begin{frame}{Qualité d'un ESG : Le Test de Martingale}
\begin{block}{Le concept du Test de Martingale}
En univers Risque Neutre, le prix actualisé de tout actif doit être égal à son prix initial (en moyenne). C'est une condition nécessaire d'absence d'arbitrage.
\end{block}

\vspace{0.3cm}
\textbf{Le Test "1 pour 1" (Zero Coupon Test) :}
Pour valider la diffusion des taux, on vérifie que la moyenne des facteurs d'actualisation stochastiques restitue la courbe des prix ZC initiale :
\[
P(0, T)_{\text{Marché}} \approx \frac{1}{N} \sum_{i=1}^{N} \exp\left(-\sum_{t=0}^{T-1} r_t^{(i)} \Delta t\right)
\]

\vspace{0.3cm}
\begin{exampleblock}{En pratique lors du TD}
Si après 1000 scénarios, votre prix calculé s'éloigne de plus de 1\% ou 2\% du prix théorique, votre ESG est :
\begin{itemize}
    \item Soit mal calibré (paramètres erronés)
    \item Soit le nombre de scénarios est insuffisant (variance trop forte)
\end{itemize}
\end{exampleblock}
\end{frame}
% ============================

% ============================================================
% SECTION 3 : VALORISATION STOCHASTIQUE
% ============================================================
\section{III. Valorisation Stochastique du BE}

\begin{frame}{Principe de Valorisation Stochastique}
\begin{definition}[BE Stochastique]
\[
\BE_{\text{stoch}} = \frac{1}{N_{\text{scen}}}\sum_{i=1}^{N_{\text{scen}}} \BE_i
\]
où $\BE_i$ est le Best Estimate calculé dans le scénario économique $i$.
\end{definition}

\vspace{0.3cm}
\textbf{Étapes de calcul :}
\begin{enumerate}
    \item \textbf{Générer $N$ scénarios} économiques (taux, actions) via ESG
    \item Pour chaque scénario $i$ :
    \begin{itemize}
        \item Projeter les flux futurs (prestations, primes, frais)
        \item Appliquer les management actions conditionnelles
        \item Actualiser les flux avec la courbe de taux du scénario $i$
        \item Obtenir $\BE_i$
    \end{itemize}
    \item \textbf{Moyenne} : $\BE_{\text{stoch}} = \bar{\BE}$
\end{enumerate}
\end{frame}

\begin{frame}{Exemple : Contrat Euro avec Garantie Plancher}
\textbf{Caractéristiques du contrat :}
\begin{itemize}
    \item PM initiale : 100 k€
    \item Taux Minimum Garanti (TMG) : 1,5\% par an
    \item Durée : 10 ans
    \item Participation aux bénéfices (PB) : 85\% des rendements au-dessus du TMG
\end{itemize}

\vspace{0.3cm}
\textbf{Actifs sous-jacents :}
\begin{itemize}
    \item 60\% obligations (rendement = taux sans risque $r_t$)
    \item 40\% actions (rendement stochastique $\mu + \sigma W_t$)
\end{itemize}

\vspace{0.3cm}
\begin{alertblock}{Question}
Quelle est la valeur du Best Estimate en tenant compte de la garantie plancher ?
\end{alertblock}
\end{frame}

\begin{frame}{Projection dans un Scénario Donné}
\textbf{Scénario $i$ : évolution du rendement actifs $R^{(i)}_t$}

\begin{table}
\centering
\scriptsize
\begin{tabular}{lccccc}
\toprule
\textbf{Année} & \textbf{PM début} & \textbf{$R^{(i)}_t$} & \textbf{TMG} & \textbf{Reval. assurée} & \textbf{PM fin} \\
\midrule
1 & 100 & 0,5\% & 1,5\% & 1,5\% & 101,5 \\
2 & 101,5 & 3,2\% & 1,5\% & 2,95\% & 104,5 \\
3 & 104,5 & -1\% & 1,5\% & 1,5\% & 106,1 \\
... & ... & ... & ... & ... & ... \\
10 & ... & 2\% & 1,5\% & 1,93\% & 130,2 \\
\bottomrule
\end{tabular}
\end{table}

\vspace{0.2cm}
\textbf{Calculs :}
\begin{itemize}
    \item Si $R^{(i)}_t < \text{TMG}$ : revalorisation = TMG (garantie activée)
    \item Si $R^{(i)}_t > \text{TMG}$ : revalorisation = TMG $+ 85\% \times (R^{(i)}_t - \text{TMG})$
\end{itemize}

\vspace{0.2cm}
Flux sortant année 10 = 130,2 k€ $\Rightarrow$ actualisé :
\[
\BE_i = \frac{130,2}{(1+r_{10})^{10}}
\]
\end{frame}

\begin{frame}{Distribution du BE sur $N$ Scénarios}
\begin{center}
\begin{tikzpicture}[scale=0.8]
    % Histogramme stylisé
    \draw[->] (0,0) -- (8,0) node[right] {$\BE_i$ (k€)};
    \draw[->] (0,0) -- (0,4) node[above] {Fréquence};
    
    % Barres
    \fill[blue!30] (1,0) rectangle (1.5,0.8);
    \fill[blue!30] (1.5,0) rectangle (2,1.5);
    \fill[blue!30] (2,0) rectangle (2.5,2.5);
    \fill[blue!30] (2.5,0) rectangle (3,3.2);
    \fill[blue!30] (3,0) rectangle (3.5,3.5);
    \fill[blue!30] (3.5,0) rectangle (4,3.2);
    \fill[blue!30] (4,0) rectangle (4.5,2.5);
    \fill[blue!30] (4.5,0) rectangle (5,1.5);
    \fill[blue!30] (5,0) rectangle (5.5,0.8);
    
    % Moyenne
    \draw[red, thick, dashed] (3.25,0) -- (3.25,3.8) node[above] {$\E[\BE]$};
    % Annotations
    \node at (1.25,-0.5) {\scriptsize 95};
    \node at (3.25,-0.5) {\scriptsize 105};
    \node at (5.25,-0.5) {\scriptsize 115};
\end{tikzpicture}
\end{center}

\vspace{0.3cm}
\begin{exampleblock}{Interprétation}
\begin{itemize}
    \item Distribution étalée : impact de la variabilité des rendements actifs
    \item \textbf{Moyenne} = BE stochastique (ex: 105 k€)
    \item Queue de distribution gauche : scénarios défavorables (garantie coûteuse)
\end{itemize}
\end{exampleblock}
\end{frame}

\begin{frame}{Convergence Monte-Carlo : Nombre de Scénarios}
\begin{block}{Erreur de Monte-Carlo}
L'erreur standard de l'estimateur est :
\[
\text{Erreur} \approx \frac{\sigma_{\BE}}{\sqrt{N_{\text{scen}}}}
\]
où $\sigma_{\BE}$ est l'écart-type des $\BE_i$.
\end{block}

\textbf{Conséquence :}
\begin{itemize}
    \item Pour diviser l'erreur par 2 : multiplier $N$ par 4
    \item Pour diviser l'erreur par 10 : multiplier $N$ par 100
\end{itemize}

\vspace{0.3cm}
\begin{table}
\centering
\begin{tabular}{lcc}
\toprule
\textbf{$N_{\text{scen}}$} & \textbf{Erreur relative} & \textbf{Usage} \\
\midrule
100 & $\sim$ 10\% & Tests rapides \\
1 000 & $\sim$ 3\% & Développement \\
\rowcolor{green!20}
5 000 - 10 000 & $<$ 1\% & \textbf{Production S2} \\
\bottomrule
\end{tabular}
\end{table}
\end{frame}

\begin{frame}{Management Actions : Décisions Conditionnelles}
\textbf{Exemples de management actions :}

\begin{enumerate}
    \item \textbf{Allocation dynamique Actif-Passif :}
    \begin{itemize}
        \item Si rendement actifs $<$ TMG : vendre actions, acheter obligations sûres
        \item Si contexte favorable : augmenter exposition actions
    \end{itemize}
    
    \item \textbf{Taux de Participation aux Bénéfices (PB) :}
    \begin{itemize}
        \item Si réserves PPB élevées : distribuer PB généreuse
        \item Si réserves faibles : PB minimale
    \end{itemize}
    
    \item \textbf{Gestion des rachats :}
    \begin{itemize}
        \item Offrir bonus de fidélité si risque de rachats massifs
    \end{itemize}
\end{enumerate}

\vspace{0.3cm}
\begin{alertblock}{Implémentation}
Les management actions doivent être \textbf{codées explicitement} dans la projection de chaque scénario :
\begin{itemize}
    \item Règles "si-alors" conditionnelles
    \item Tests de variables d'état (réserves, rendements cumulés, ...)
\end{itemize}
\end{alertblock}
\end{frame}

\begin{frame}[fragile]{Pseudo-code : Projection avec Management Actions}
\begin{lstlisting}[style=pythonstyle, basicstyle=\ttfamily\tiny]
for scenario in range(N_scenarios):
    # Initialisation
    PM = 100  # Provision Mathematique
    PPB = 0   # Provision PB differee
    
    for annee in range(1, duree+1):
        # 1. Calculer rendement actifs (depend du scenario)
        rend_actifs = calc_rendement(scenario, annee)
        
        # 2. Management Action : allocation dynamique
        if rend_actifs < TMG:
            pct_actions = 20%  # Reduction exposition actions
        else:
            pct_actions = 40%  # Normal
        
        # 3. Calcul revalorisation contrat
        if rend_actifs < TMG:
            reval_contrat = TMG  # Garantie activee
            cout_garantie = PM * (TMG - rend_actifs)
            PPB -= cout_garantie  # Prelevement sur PPB
        else:
            surplus = rend_actifs - TMG
            PB_distribuee = 0.85 * surplus
            reval_contrat = TMG + PB_distribuee
            PPB += 0.15 * surplus  # Dotation PPB
        
        # 4. Mise a jour PM
        PM = PM * (1 + reval_contrat)
    
    # Actualisation flux final
    BE_scenario[scenario] = PM / (1 + taux_actualisation)**duree

# BE stochastique = moyenne
BE_stoch = np.mean(BE_scenario)
\end{lstlisting}
\end{frame}

% === NOUVELLE SLIDE (3/5) ===
\begin{frame}{Défis de Production : Nested Simulations vs Proxies}
\textbf{Le problème du calcul du SCR :}
Le SCR est une VaR à 1 an des Fonds Propres (NAV).
\begin{itemize}
    \item Pour calculer la NAV en $t=1$ (dans un scénario de choc), il faut recalculer le BE.
    \item Recalculer le BE implique de relancer 1000 scénarios (de $t=1$ à $t=T$).
\end{itemize}

\vspace{0.2cm}
\textbf{Complexité "Stochastic in Stochastic" (Nested) :}
\[ \text{Complexité} = N_{\text{scénarios SCR}} \times N_{\text{scénarios BE}} \approx 100\,000 \times 1\,000 = 10^8 \text{ projections} \]
$\Rightarrow$ \alert{Impossible} en temps de calcul raisonnable.

\vspace{0.3cm}
\begin{block}{La solution industrielle : Les Proxies (Metamodèles)}
On remplace le moteur de projection par une fonction mathématique $f$ rapide à évaluer :
\[ \text{NAV}(X_{\text{eco}}) \approx f(\text{Taux}, \text{Actions}, \text{Volatilité}...) \]
Techniques : \textbf{Least Squares Monte Carlo (LSMC)}, \textbf{Polynômes de Réplication}.
\end{block}
\end{frame}
% ============================

\begin{frame}{Limites et Améliorations}
\textbf{Défis de la valorisation stochastique :}

\begin{itemize}
    \item[\texttimes] \textbf{Temps de calcul} : projection de milliers de scénarios $\times$ millions de polices
    \item[\texttimes] \textbf{Complexité des modèles} : nombreux paramètres à calibrer
    \item[\texttimes] \textbf{Modèles imparfaits} : hypothèses simplificatrices (ex: lois de rachats)
\end{itemize}

\vspace{0.3cm}
\textbf{Techniques d'accélération :}

\begin{enumerate}
    \item \textbf{Model Points} : regrouper polices similaires (déjà vu séance 4)
    \item \textbf{Least Squares Monte-Carlo (LSMC)} : régression pour approximer valeurs futures
    \item \textbf{Réplication de portefeuille} : approximer passif par portefeuille d'actifs répliquant
    \item \textbf{Calcul parallèle / GPU} : exploiter puissance de calcul moderne
\end{enumerate}

\vspace{0.3cm}
\begin{exampleblock}{En pratique}
Compromis entre précision et temps de calcul : valorisation détaillée trimestrielle + approximations mensuelles.
\end{exampleblock}
\end{frame}

% ============================================================
% SECTION 4 : ALM EN ASSURANCE VIE
% ============================================================
\section{IV. Asset-Liability Management (ALM)}

\begin{frame}{Qu'est-ce que l'ALM ?}
\begin{definition}[ALM - Gestion Actif-Passif]
L'\textbf{ALM} est l'ensemble des stratégies et techniques permettant de gérer conjointement les actifs (placements) et les passifs (engagements) pour :
\begin{itemize}
    \item Minimiser les risques de désadéquation
    \item Maximiser la rentabilité ajustée au risque
    \item Assurer la solvabilité à long terme
\end{itemize}
\end{definition}

\vspace{0.3cm}
\begin{block}{Contexte Assurance Vie}
\begin{itemize}
    \item \textbf{Passif} : engagements de long terme (20-40 ans) avec garanties
    \item \textbf{Actif} : portefeuille d'obligations, actions, immobilier
    \item \textbf{Objectif} : faire en sorte que les actifs "financent" correctement les passifs
\end{itemize}
\end{block}
\end{frame}

\begin{frame}{Les Risques en ALM}
\begin{columns}
\column{0.5\textwidth}
\textbf{1. Risque de Taux :}
\begin{itemize}
    \item Baisse des taux : réinvestissement à taux faibles
    \item Hausse des taux : moins-values latentes sur obligations
\end{itemize}

\textbf{2. Risque de Liquidité :}
\begin{itemize}
    \item Rachats massifs imprévus
    \item Nécessité de vendre actifs en perte
\end{itemize}

\textbf{3. Risque de Marché :}
\begin{itemize}
    \item Chute des marchés actions
    \item Impact sur capacité à servir PB
\end{itemize}

\column{0.5\textwidth}
\begin{center}
\begin{tikzpicture}[scale=0.6]
    % Schéma Actif-Passif
    \node[draw, fill=blue!20, rectangle, minimum width=3cm, minimum height=1.5cm] (actif) at (0,2) {\textbf{ACTIF}};
    \node[draw, fill=red!20, rectangle, minimum width=3cm, minimum height=1.5cm] (passif) at (0,-2) {\textbf{PASSIF}};
    
    \node[align=center] at (0,2) {\scriptsize Obligations\\\scriptsize Actions\\\scriptsize Immobilier};
    \node[align=center] at (0,-2) {\scriptsize Provisions\\\scriptsize Mathématiques\\\scriptsize + Risk Margin};
    
    % Flèches risques
    \draw[->, thick, red] (actif) -- (passif) node[midway, right] {\scriptsize Matching};
    \draw[->, thick, orange, dashed] (3,2) -- (3,-2) node[midway, right] {\scriptsize Risques};
\end{tikzpicture}
\end{center}
\end{columns}
\end{frame}

% === NOUVELLE SLIDE (4/5) ===
\begin{frame}{L'ALM comme outil d'optimisation du SCR}
\textbf{Le lien direct ALM $\leftrightarrow$ SCR :}
\begin{itemize}
    \item Le \textbf{SCR Taux} correspond à la perte maximale de Fonds Propres (NAV) en cas de choc de taux.
    \item $\text{NAV} = \text{Valeur Actif} - \text{Best Estimate}$.
\end{itemize}

\vspace{0.3cm}
\begin{block}{Mécanisme d'immunisation}
Si l'ALM est parfaite (Duration Gap $\approx 0$ et Convexité maîtrisée) :
\[ \Delta \text{Actif} \approx \Delta \text{Best Estimate} \]
\[ \Rightarrow \Delta \text{NAV} \approx 0 \]
\end{block}

\vspace{0.3cm}
\begin{alertblock}{Conséquence stratégique}
Une bonne gestion ALM ne sert pas uniquement à sécuriser les engagements, elle permet de \textbf{réduire drastiquement le besoin en capital} (SCR Market).
\end{alertblock}
\end{frame}
% ============================

\begin{frame}{Duration Gap : Concept Clé}
\begin{definition}[Duration]
La duration d'un flux de paiements est la moyenne pondérée des échéances, pondérée par les valeurs actuelles :
\[
D = \frac{\sum_{t} t \cdot \text{VA}(F_t)}{\sum_{t} \text{VA}(F_t)}
\]
\end{definition}

\vspace{0.3cm}
\begin{block}{Duration Gap}
\[
\text{Duration Gap} = D_{\text{Actif}} - D_{\text{Passif}}
\]

\textbf{Interprétation :}
\begin{itemize}
    \item Si Gap $> 0$ : actifs plus longs que passifs $\Rightarrow$ risque si taux monte
    \item Si Gap $< 0$ : actifs plus courts que passifs $\Rightarrow$ risque de réinvestissement
    \item \textbf{Objectif} : Gap $\approx 0$ pour immunisation
\end{itemize}
\end{block}
\end{frame}

\begin{frame}{Exemple Chiffré : Duration Gap}
\textbf{Passif : Rente viagère}
\begin{itemize}
    \item Paiements annuels de 10 k€ sur 20 ans
    \item Taux d'actualisation 2\%
    \item Calcul $\Rightarrow$ $D_{\text{Passif}} \approx 12$ ans
\end{itemize}

\vspace{0.3cm}
\textbf{Actif : Obligations}
\begin{itemize}
    \item Portefeuille d'obligations de maturités variées
    \item Duration moyenne du portefeuille : $D_{\text{Actif}} = 8$ ans
\end{itemize}

\vspace{0.3cm}
\begin{alertblock}{Calcul du Gap}
\[
\text{Gap} = 8 - 12 = -4 \text{ ans}
\]
$\Rightarrow$ Actifs trop courts : \textbf{risque de réinvestissement} si taux baisse.
\end{alertblock}

\textbf{Action ALM :} Acheter obligations plus longues (maturité 15-20 ans) pour allonger $D_{\text{Actif}}$.
\end{frame}

\begin{frame}{Stratégies d'Immunisation}
\textbf{Objectif : } Rendre le bilan insensible aux variations de taux.
\begin{block}{Immunisation de Redington (1952)}
Conditions pour immunisation parfaite :
\begin{enumerate}
    \item \textbf{Valeurs actuelles égales} : $\text{VA}_{\text{Actif}} = \text{VA}_{\text{Passif}}$
    \item \textbf{Durations égales} : $D_{\text{Actif}} = D_{\text{Passif}}$
    \item \textbf{Convexité supérieure} : $C_{\text{Actif}} > C_{\text{Passif}}$
\end{enumerate}
\end{block}

\vspace{0.3cm}
\textbf{Explication intuitive :}
\begin{itemize}
    \item Condition 1 : adéquation initiale
    \item Condition 2 : sensibilité identique aux petits chocs de taux
    \item Condition 3 : protection supplémentaire contre grands chocs (convexité = courbure)
\end{itemize}

\vspace{0.3cm}
\begin{exampleblock}{Limite}
Immunisation parfaite difficile en pratique : nécessite rééquilibrage fréquent du portefeuille.
\end{exampleblock}
\end{frame}

\begin{frame}{Convexité : Concept Avancé}
\begin{definition}[Convexité]
Mesure la courbure de la relation prix-taux :
\[
C = \frac{1}{P}\frac{\partial^2 P}{\partial r^2}
\]
\end{definition}

\textbf{Impact :}
\begin{itemize}
    \item Convexité élevée $\Rightarrow$ gains asymétriques
    \item Si taux baisse : prix monte beaucoup (bon)
    \item Si taux monte : prix baisse peu (protection)
\end{itemize}

\vspace{0.3cm}
\begin{center}
\begin{tikzpicture}[scale=0.8]
    \draw[->] (0,0) -- (5,0) node[right] {Taux};
    \draw[->] (0,0) -- (0,4) node[above] {Prix};
    
    % Courbe convexe
    \draw[thick, blue] (0.5,3.5) .. controls (2.5,2) .. (4.5,1);
    % Tangente
    \draw[dashed, red] (0,4) -- (5,0.5);
    
    \node[blue] at (1,3) {Convexe};
    \node[red] at (3.5,0.3) {Linéaire};
\end{tikzpicture}
\end{center}
\end{frame}

\begin{frame}{Cash-Flow Matching vs Duration Matching}
\begin{columns}
\column{0.5\textwidth}
\textbf{Cash-Flow Matching :}
\begin{itemize}
    \item Faire correspondre \textbf{exactement} chaque flux d'actif à chaque flux de passif
    \item Très contraignant
    \item Immunisation parfaite
\end{itemize}

\textbf{Avantages :}
\begin{itemize}
    \item[\checkmark] Risque minimal
\end{itemize}

\textbf{Inconvénients :}
\begin{itemize}
    \item[\texttimes] Peu flexible
    \item[\texttimes] Rendement sous-optimal
\end{itemize}

\column{0.5\textwidth}
\textbf{Duration Matching :}
\begin{itemize}
    \item Égaliser les \textbf{durations} seulement
    \item Plus flexible
    \item Immunisation approximative
\end{itemize}

\textbf{Avantages :}
\begin{itemize}
    \item[\checkmark] Meilleur rendement
    \item[\checkmark] Gestion simplifiée
\end{itemize}

\textbf{Inconvénients :}
\begin{itemize}
    \item[\texttimes] Risque résiduel
    \item[\texttimes] Nécessite rééquilibrage
\end{itemize}
\end{columns}

\vspace{0.3cm}
\begin{exampleblock}{En pratique}
Compromis hybride : duration matching avec contraintes de liquidité (cash-flow partiel).
\end{exampleblock}
\end{frame}

\begin{frame}{Matching Adjustment (MA) et Volatility Adjustment (VA)}
\textbf{Mécanismes Solvency II pour stabiliser le bilan :}

\begin{block}{Matching Adjustment (MA)}
\begin{itemize}
    \item Permet d'actualiser le passif avec un taux \textbf{supérieur} au taux sans risque
    \item Applicable si \textbf{adossement actif-passif très strict} (cash-flow matching)
    \item Réduit artificiellement le BE $\Rightarrow$ améliore ratio de solvabilité
\end{itemize}
Formule :
\[
r_{\text{actualisation}} = r_{\text{EIOPA}} + \text{MA}
\]
où MA = spread actifs - spread de risque de crédit.
\end{block}

\vspace{0.3cm}
\begin{block}{Volatility Adjustment (VA)}
\begin{itemize}
    \item Ajustement symétrique de la courbe de taux pour absorber volatilité excessive
    \item Automatique (pas de conditions strictes comme MA)
    \item Effet plus faible que MA
\end{itemize}
\end{block}
\end{frame}

\begin{frame}{Avantages et Critiques de MA/VA}
\begin{columns}
\column{0.5\textwidth}
\textbf{Avantages :}
\begin{itemize}
    \item[\checkmark] Stabilité du bilan face aux fluctuations de spread
    \item[\checkmark] Évite pro-cyclicité
    \item[\checkmark] Réduit artificialité de la valorisation économique pure
\end{itemize}

\column{0.5\textwidth}
\textbf{Critiques :}
\begin{itemize}
    \item[\texttimes] Masque une partie du risque réel
    \item[\texttimes] Complexité réglementaire
    \item[\texttimes] Asymétrie : bénéfice en période de stress, mais pas en période favorable
\end{itemize}
\end{columns}

\vspace{0.3cm}
\begin{alertblock}{Réforme Solvency II 2025-2026}
\begin{itemize}
    \item Révision des règles de MA (assouplissement conditions)
    \item Ajustements de VA pour mieux refléter conditions de marché
    \item Objectif : encourager investissements long terme (infrastructures, transition énergétique)
\end{itemize}
\end{alertblock}
\end{frame}

\begin{frame}{ALM Dynamique : Approche Stochastique}
\textbf{ALM classique : } Statique, basée sur durations instantanées.
\textbf{ALM dynamique : } Intègre évolution stochastique des variables économiques.

\vspace{0.3cm}
\begin{block}{Principe}
\begin{enumerate}
    \item Simuler $N$ scénarios économiques (taux, actions)
    \item Pour chaque scénario :
    \begin{itemize}
        \item Projeter actif et passif simultanément
        \item Appliquer règles de gestion d'actifs conditionnelles (rééquilibrage)
        \item Calculer ratio de couverture, SCR, ...
    \end{itemize}
    \item Analyser distribution des résultats :
    \begin{itemize}
        \item VaR du ratio de couverture
        \item Probabilité de déficit
    \end{itemize}
\end{enumerate}
\end{block}

\vspace{0.3cm}
\textbf{Avantage : } Capture interactions complexes entre actifs, passifs et décisions de gestion.
\end{frame}

\begin{frame}{Indicateurs ALM Clés}
\begin{table}
\centering
\small
\begin{tabular}{lp{6cm}c}
\toprule
\textbf{Indicateur} & \textbf{Définition} & \textbf{Cible} \\
\midrule
Duration Gap & $D_A - D_P$ & $\approx 0$ \\
\midrule
Ratio de couverture & $\frac{\text{Valeur Actifs}}{\text{Valeur Passifs}}$ & $> 1$ \\
\midrule
ALM Gap (VaR) & VaR 99,5\% de $(A - P)$ & Minimiser \\
\midrule
Liquidité & \% actifs liquides sur total & $> 20\%$ \\
\midrule
Taux de recouvrement & Capacité à honorer engagements en stress & $100\%$ \\
\bottomrule
\end{tabular}
\end{table}

\vspace{0.3cm}
\begin{exampleblock}{Dashboard ALM}
Les assureurs utilisent des tableaux de bord temps réel avec ces indicateurs pour piloter l'ALM.
\end{exampleblock}
\end{frame}

\begin{frame}{Cas Pratique : Étude ALM Simplifiée}
\textbf{Contexte :}
\begin{itemize}
    \item Assureur Vie avec BE = 1 000 M€, Duration passif = 15 ans
    \item Actifs : 1 050 M€, Duration actif = 10 ans
    \item Stress taux : +100 bp (hausse 1\%)
\end{itemize}

\vspace{0.3cm}
\textbf{Impact sur Actifs :}
\[
\Delta A \approx -D_A \times \Delta r \times A = -10 \times 0.01 \times 1050 = -105 \text{ M€}
\]
$\Rightarrow$ Actifs = 945 M€

\vspace{0.3cm}
\textbf{Impact sur Passifs :}
\[
\Delta P \approx -D_P \times \Delta r \times P = -15 \times 0.01 \times 1000 = -150 \text{ M€}
\]
$\Rightarrow$ Passifs = 850 M€

\vspace{0.3cm}
\begin{alertblock}{Conclusion}
Ratio couverture passe de $\frac{1050}{1000} = 1.05$ à $\frac{945}{850} = 1.11$. \textbf{Amélioration !}

$\Rightarrow$ Duration gap négatif protège en cas de hausse de taux.
\end{alertblock}
\end{frame}

% ============================================================
% SECTION 5 : SYNTHÈSE ET EXERCICES
% ============================================================
\section{Synthèse et Travaux Dirigés}

\begin{frame}{Synthèse de la Séance}
\textbf{Concepts clés abordés :}

\begin{enumerate}
    \item \textbf{Limites déterministes} : insuffisance pour options, garanties, management actions
    \item \textbf{ESG} : Calibration Risque Neutre, Test de Martingale
    \item \textbf{Monte-Carlo} : simulation de scénarios multiples pour valorisation
    \item \textbf{BE stochastique} : moyenne des BE sur tous scénarios
    \item \textbf{ALM} : gestion actif-passif, duration gap, immunisation et lien SCR
    \item \textbf{Proxies} : LSMC pour accélérer les calculs
\end{enumerate}

\vspace{0.3cm}
\begin{block}{Compétences acquises}
\begin{itemize}
    \item Comprendre la nécessité d'approche stochastique
    \item Implémenter simulations Monte-Carlo en Python
    \item Analyser risques ALM et proposer stratégies d'immunisation
\end{itemize}
\end{block}
\end{frame}

\begin{frame}{Transition vers les TD}
\textbf{Exemples de Travaux Dirigés :}

\begin{block}{Exercice 5.1 : Simulations Monte-Carlo - Garantie Plancher }
\begin{itemize}
    \item Implémenter ESG Hull-White en Python (avec test de martingale !)
    \item Simuler 1 000 scénarios de taux
    \item Projeter contrat euro avec TMG
    \item Calculer BE stochastique et comparer à approche déterministe
\end{itemize}
\end{block}

\begin{block}{Exercice 5.2 : Management Actions et PB Dynamique }
\begin{itemize}
    \item Modéliser règles de PB conditionnelles
    \item Intégrer Provision pour PB Différée (PPB)
    \item Analyser impact sur distribution du BE
\end{itemize}
\end{block}
\end{frame}

\begin{frame}{Vers les TD}
\begin{block}{Exercice 5.3 : Duration Gap et Stratégie ALM (35 min)}
\begin{itemize}
    \item Calculer durations actif et passif
    \item Simuler impact de chocs de taux
    \item Proposer allocation optimale pour immunisation
\end{itemize}
\end{block}


\vspace{0.5cm}
vous pouvez vous inspirer du notebook proposé.
\end{frame}

\begin{frame}{Conseils Pratiques pour les TD}
\textbf{Pour réussir les exercices :}

\begin{itemize}
    \item \textbf{Bien lire les énoncés} : hypothèses, paramètres donnés
    \item \textbf{Décomposer les problèmes} : étapes simples plutôt qu'un gros code
    \item \textbf{Vérifier la cohérence} : utiliser le Test de Martingale
    \item \textbf{Performance Python} : Utiliser la vectorisation
    \item \textbf{Visualiser les résultats} : matplotlib pour comprendre distributions
\end{itemize}

\vspace{0.3cm}
\begin{alertblock}{Piège courant}
Oublier d'actualiser les flux futurs $\Rightarrow$ toujours diviser par $(1+r)^t$ !
\end{alertblock}

\vspace{0.3cm}
% \begin{exampleblock}{Astuce Python}
% \texttt{np.random.seed(42)} au début pour reproductibilité des simulations aléatoires.
% \end{exampleblock}
\end{frame}

% === NOUVELLE SLIDE (5/5) ===
\begin{frame}[fragile]{Bonnes Pratiques de Code : La Vectorisation}
\textbf{Pourquoi vectoriser ?}
Python est lent sur les boucles \texttt{for}. NumPy est optimisé en C.
Le temps de calcul peut être divisé par 100 !

\begin{columns}
\column{0.5\textwidth}
\textbf{\textcolor{red}{Mauvaise pratique (Lent)}}
\begin{lstlisting}[style=pythonstyle]
# Générer 1000 scénarios
R = []
for i in range(1000):
    scenario = []
    for t in range(40):
        val = generer_un_pas()
        scenario.append(val)
    R.append(scenario)
\end{lstlisting}

\column{0.5\textwidth}
\textbf{\textcolor{green}{Bonne pratique (Rapide)}}
\begin{lstlisting}[style=pythonstyle]
# Générer tout d'un coup
# (N_scen, N_steps)
Z = np.random.normal(0, 1, (1000, 40))
# Calcul vectoriel
R = r0 + drift * dt + sigma * Z
\end{lstlisting}
\end{columns}

\vspace{0.3cm}
\textit{En entretien technique (assurance/banque), la non-vectorisation est souvent éliminatoire.}
\end{frame}
% ============================

\begin{frame}{Pour Aller Plus Loin}
\textbf{Lectures recommandées :}

\begin{enumerate}
    \item \textit{Interest Rate Models - Theory and Practice} (Brigo \& Mercurio)
    \item \textit{Stochastic Calculus for Finance II} (Shreve)
    \item Articles EIOPA : "Guidelines on use of ESG"
    \item IA (Institut des Actuaires) - mémoires sur ESG et ALM
    \item Borel-Mathurin F., Vedani J., "Market-consistent valuation: a step towards calculation stability" \texttt{https://hal.science/hal-02282378}
    \item Borel-Mathurin F., El Karoui N., Loisel S. et Vedani J. (2020), « Locality in Time of the European
Insurance Regulation Risk-Neutral Valuation Framework, a pre- and post-Covid Analysis and Further Developments »
\end{enumerate}

\vspace{0.3cm}
\textbf{Logiciels professionnels :}
\begin{itemize}
    \item \textbf{Prophet} (FIS) : modélisation actif-passif
    \item \textbf{MoSes} (Moody's) : ESG et ALM
    \item \textbf{AAA} (Milliman) : valorisation stochastique
\end{itemize}

% \vspace{0.3cm}
% \textbf{Certifications :}
% \begin{itemize}
%     \item CFA (Chartered Financial Analyst) - niveau II/III couvre ALM
%     \item IA (Institut des Actuaires) - mémoires sur ESG et ALM
% \end{itemize}
\end{frame}

\begin{frame}{Questions \& Discussion}
\begin{center}
\Huge Merci pour votre attention !
\vspace{1cm}

\Large Questions ?

\vspace{1cm}

\normalsize
\textit{Prochaine séance : Séance 6 - Risk Margin et Provisions Consolidées}

\vspace{0.5cm}

Contact : \texttt{fabrice.mathurin@gmail.com}
\end{center}
\end{frame}

% ============================================================
% ANNEXES (optionnel)
% ============================================================
\appendix

\begin{frame}{Annexe A : Formules Mathématiques Détaillées}
\textbf{Solution analytique Vasicek :}
\[
r_t = r_0 e^{-at} + b(1 - e^{-at}) + \sigma \int_0^t e^{-a(t-s)} dW_s
\]

\vspace{0.3cm}
\textbf{Prix ZC-Bond dans Vasicek :}
\[
P(t, T) = A(t,T) e^{-B(t,T) r_t}
\]
avec :
\[
B(t,T) = \frac{1 - e^{-a(T-t)}}{a}, \quad A(t,T) = \exp\left[\frac{(B(t,T) - T + t)(a^2 b - \sigma^2/2)}{a^2} - \frac{\sigma^2 B(t,T)^2}{4a}\right]
\]

\vspace{0.3cm}
\textbf{Variance asymptotique CIR :}
\[
\Var(r_\infty) = \frac{b\sigma^2}{2a}
\]
\end{frame}

\begin{frame}{Annexe B : Code Python - ESG Complet}
\textbf{Exemple de structure pour ESG multi-variables :}

\begin{itemize}
    \item Classe \texttt{ESG} avec méthodes :
    \begin{itemize}
        \item \texttt{calibrate()}: calibration paramètres sur marché
        \item \texttt{generate\_scenarios(N)}: génération $N$ scénarios (vectorisé)
        \item \texttt{martingale\_test()}: validation auto
        \item \texttt{get\_interest\_rates()}: extraction taux
        \item \texttt{get\_equity\_returns()}: extraction rendements actions
    \end{itemize}
    \item Gestion corrélations via matrice de Cholesky
    \item Export résultats en DataFrame pandas
\end{itemize}

\vspace{0.3cm}
\textit{Exemples de code disponibles dans le notebook Seance5\_Codes\_Python.ipynb}
\end{frame}

\begin{frame}{Annexe C : Glossaire}
\begin{description}
    \item[ALM] Asset-Liability Management (Gestion Actif-Passif)
    \item[BE] Best Estimate
    \item[CIR] Cox-Ingersoll-Ross (modèle de taux)
    \item[ESG] Economic Scenario Generator
    \item[MA] Matching Adjustment
    \item[PB] Participation aux Bénéfices
    \item[PPB] Provision pour Participation aux Bénéfices (différée)
    \item[TMG] Taux Minimum Garanti
    \item[VA] Volatility Adjustment
    \item[VaR] Value at Risk
\end{description}
\end{frame}

\end{document}
```